% !TEX root =  master.tex
\chapter{Systemarchitektur} \label{sec:system}

Für das System musste eine Gesamte Architektur entwickelt werden, welche die Komponenten HB10, Nutzer (Webseite) und Server miteinander verbindet (s. \ref{fig:architektur}). Die einzige bereits vorhandene Komponente ist der HB10, welcher allerdings noch keine Kamera enthält. 
In diesem Kapitel wird ein grober Überblick über die Architektur gegeben, die Details zu den einzelnen Komponenten werden zu einem späteren Zeitpunkt erläutert. 

\begin{figure}[H]
    \centering
    \includesvg[width=\textwidth]{img/System.svg}
    \caption[Gesamtarchitektur der Video-on-Demand-Lösung]{Gesamtarchitektur der entwickelten Video-on-Demand-Lösung. Der HB10 erfasst und überträgt die aufgenommenen Videodaten an einen Spring-Boot-Server, der die Videos im Dateisystem und die zugehörigen Metadaten in einer Datenbank speichert. 
    Nutzer greifen über einen Browser auf die Webplattform zu und können die bereitgestellten Videos über den Server abrufen.}
    \label{fig:architektur}
\end{figure}

Der HB10 kann seine aufgenommenen Videos mittels HTTP an den Server übertragen. Als Server dient ein einfacher Webserver, welcher Zugriff auf eine Datenbank und das Dateisystem hat, um empfangene Videos und deren Metadaten zu speichern. Der Nutzer erhält über den Webserver Zugriff auf eine Webplattform, 
über welche er die Videos ansehen kann.

Der HB10 ist dafür zuständig, den Spielbetrieb durchzuführen, und die Videos für die Dart-Würfe aufzunehmen, diese zu komprimieren und an den Server zu senden. Der Server speichert die Videos inklusive der zugehörigen Metadaten und stellt sie den Nutzern des HB10 über eine Webseite bereit. 
Der Nutzer kann auf dieser Webseite die Metadaten abrufen, ein gewünschtes Video auffinden und anschließend ansehen.

\section{Nutzungsszenarien}

Das zentrale Nutzungsszenario besteht in der nachträglichen Überprüfung eines asynchron durchgeführten Turnierspiels. Der Spieler führt das Spiel am HB10 durch, während die relevanten Spielsituationen durch Videoaufnahmen dokumentiert werden. 
Die technische Verarbeitung der einzelnen Aufnahmen sowie deren Übertragung und Speicherung erfolgt entsprechend der zuvor beschriebenen Systemarchitektur. 

Nach der Bereitstellung des Videos können der aufgezeichnete Spieler und weitere für das Turnier berechtigte Teilnehmer die Aufnahme über die Webplattform aufrufen. Dort kann der dokumentierte Spielverlauf betrachtet und mit dem gemeldeten Spielergebnis verglichen werden. 
Der Zugriff auf die Aufnahme dient dabei der manuellen Prüfung, ob aus Sicht des Nutzers Hinweise auf einen möglichen Betrugsversuch vorliegen. 

Wird eine Spielsituation als auffällig eingeschätzt, kann das betroffene Spiel zur weiteren Überprüfung gemeldet werden. Eine solche Meldung stellt noch keine abschließende Bewertung dar, sondern leitet den Vorgang an eine zuständige Prüfinstanz weiter. 
Diese kann das gespeicherte Videomaterial ebenfalls einsehen und anschließend entscheiden, ob das gemeldete Spiel als legitim anzusehen ist. 

An dem Nutzungsszenario sind somit drei Nutzergruppen beteiligt. Der Spieler erzeugt durch die Teilnahme am Turnierspiel die aufzunehmenden Spieldaten. Die berechtigten Teilnehmer verwenden die Webplattform zur Betrachtung der bereitgestellten Aufnahme und können auffällige Situationen melden. 
Die zuständige Prüfinstanz übernimmt schließlich die abschließende Bewertung des gemeldeten Spiels. Das System stellt hierfür das Videomaterial und die zugehörigen Spielinformationen bereit, nimmt jedoch keine automatische Entscheidung über das Vorliegen eines Betrugsversuches vor.

\newpage